\section{Why PPO Falls Behind}
\label{sec:analysis}

\paragraph{Where the return goes.}
Splitting each episode's return into network utility (delivery, disconnection,
SLA, utilization and overload terms), shaping, and move costs
(Table~\ref{tab:decomp}) localizes PPO's deficit. Averaged over roots, the
bandit earns a utility of \UtilBandit{} and pays \CostBandit{} in move and
reversal costs; PPO earns \UtilPPO{} and pays \CostPPO. Of the
\Gap-point gap, \GapFromCosts{} (\GapCostSharePct\,\%) is move and reversal
cost and \GapFromUtility{} is network utility. PPO's utility varies widely
across roots (\PPOUtilMin{} to \PPOUtilMax, against \BanditUtilMin{} to
\BanditUtilMax{} for the bandit); its best root is above the bandit's best. On \PPOFlappingRoots{} of \TestPPORoots{} roots more
than half of its accepted moves are \emph{reversals} (up to \PPORevShareMax\,\%),
and its per-root move costs range from \PPOCostMin{} to \PPOCostMax{}, against
\BanditCostMin{} to \BanditCostMax{} for the bandit. The shaping term
contributes at most \ShapingMaxAbs{} per episode to the evaluated return of
any policy, but that does not bound its effect on learning: a learner trained
on the shaped reward sees it at every step, and removing it changes the
bandit's return by \RQThreeNoShaping{} on two roots, within learner noise
(\S\ref{sec:rq3}).

\begin{table}[t]
\centering\small
\caption{Return decomposition on the test episodes (means over episodes and,
for learners, over roots). Utility = delivery + disconnection + SLA +
utilization + overload terms; Costs = move and reversal terms.}
\label{tab:decomp}
\setlength{\tabcolsep}{4pt}
\begin{adjustbox}{max width=\columnwidth}
\begin{tabular}{@{}lrrrrrr@{}}
\toprule
Policy & Return & Utility & Costs & Shaping & Moves & Reversals \\
\midrule
\tabinput{tables/decomposition}
\bottomrule
\end{tabular}
\end{adjustbox}
\end{table}

\paragraph{A degenerate oscillation.}
Figure~\ref{fig:case} shows one test episode of the deceptive-local-optimum
scenario. After moving at nearly every opportunity in the first hour, PPO
settles into flipping one demand between two candidates every three
intervals---exactly when the hold-down expires---and pays the reversal price
each time. The bandit makes a handful of moves at the burst onset and
otherwise holds its configuration. The behaviour is not an artefact of checkpoint selection:
\PPOFinalFlappingRoots{} of \PPOFinalRoots{} final (400k) PPO checkpoints also
reverse more than half of their moves.

\paragraph{A hypothesis: small immediate costs in noisy advantages.}
A move that returns a demand to its previous path costs at least $0.38$
reward units once (fixed and reversal terms of $C^{\mathrm{TE}}$,
Eq.~\ref{eq:reward}), plus its traffic and divergence terms. PPO's advantage
estimates aggregate, through GAE with $\gamma\lambda\approx0.945$, the
fluctuations of a reward whose per-interval mean alone ranges from
\PerIntervalMin{} to \PerIntervalMax{} across test episodes, driven by bursts,
failures and diurnal load. We hypothesise that a one-off cost of this size is
small relative to the variance of the advantage target, so that the policy
gradient does not attribute it to the step that incurs it. The bandit's
regression target \emph{is} that step's reward, so the cost enters its
estimate without noise from future intervals. We did not measure advantage
noise directly; \S\ref{sec:rq2} tests the explanation indirectly by
removing PPO's horizon, and \S\ref{sec:discussion} discusses what speaks
against it.

\paragraph{What the learners' actions are worth.}
On states visited by each learner on test episodes we compute, by cloning the
simulator, the exact one-interval reward of every legal action and the
six-interval value of the learner's action, of the exact myopic-best action,
and of no-op (Table~\ref{tab:fidelity}). Neither learner is close to the exact myopic
optimum, but they fail in opposite directions. The bandit is conservative: it
chooses no-op in \FidBanditNoopPolicy\,\% of states although no-op is optimal
in only \FidBanditNoopOptimal\,\%, and it matches the optimum in
\FidBanditTopOne\,\% of states (one-interval regret \FidBanditRegret). PPO
over-acts: it chooses no-op in \FidPPONoopPolicy\,\% of states, matches the
optimum in \FidPPOTopOne\,\% (regret \FidPPORegret), and over six intervals
its chosen actions are worth \emph{less than doing nothing} on
\FidPPONegRoots{} of \FidPPORoots{} roots (mean $\Delta G_6=\FidPPOGSix$
against \FidBanditGSix{} for the bandit and \FidPPOGSixMyopic{} for the exact
myopic choice on PPO's states). PPO's frequent moves therefore do not buy
delayed value that a myopic criterion would miss.

\begin{table}[t]
\centering\small
\caption{What the learners' actions are worth, from exact simulator
rollouts on states each learner visits (test seeds 3001--3003, every third
step). Top-1: share of states in which the learner's action is the exact
one-interval optimum; regret: exact one-interval reward lost; no-op: share
of states in which the learner / the optimum chooses no-op; $\Delta G_6$:
six-interval value over no-op (no-op continuation) of the learner's action
and of the exact myopic-best action.}
\label{tab:fidelity}
\setlength{\tabcolsep}{3.5pt}
\begin{adjustbox}{max width=\columnwidth}
\begin{tabular}{@{}lrrrrrr@{}}
\toprule
 & Root & Top-1 & Regret & No-op (learner/opt.) & $\Delta G_6$ learner & $\Delta G_6$ myopic \\
\midrule
\tabinput{tables/model_fidelity}
\bottomrule
\end{tabular}
\end{adjustbox}
\end{table}

\begin{figure*}[t]
\centering
\includegraphics[width=\textwidth]{fig_case}
\caption{One test episode (deceptive-local-optimum, seed 3001; learners of root
314159). Top: worst link utilization; bottom: accepted moves ($\times$ = return to
the previous path within six intervals). PPO flips one demand every three
intervals, the hold-down period; the bandit makes few moves. Oracle-1 sees the
next interval's traffic and is shown for reference.}
\label{fig:case}
\end{figure*}
